\scanpage{50}

% B1 | text | two opening lines
Donc on peut \deleted{considérer que} \inserted{identifier} les $(C,\alpha)$ \deleted{\unsure{sont}} entre
eux \deleted{\unsure{à deux}}, on trouve donc

% B2 | display | upper centre
\[
[C(\xi),\alpha(\xi)]\qquad\text{canonique.}
\]

% B3 | text | boxed and crossed opening of abandoned paragraph
\deleted{On peut \unsure{donc|alors} considérer sur $\mathbb A''$\inserted{$=\mathbb A'\times A_\infty$} (anneau des adèles
premières à $p$), \deleted{\unsure{un}} la gerbe $g''$ \inserted{de groupe $\mu_2$} qui sur $A'$ est la
gerbe des racines carrées de $A'(1)$, et sur $A_\infty$ est la
gerbe type $g'_\infty$ déjà envisagée, alors}

% B4 | text | main paragraph, upper middle
En fonction de ce choix, on trouve une gerbe
canonique $\unsure{\mathscr G|\mathfrak G|\mathcal G}(\xi)\deleted{\unsure{=}}$ sur l'anneau des adèles,
de groupe $\mu_2$, savoir la gerbe qui sur $A''=A'\times A_\infty$ est
la gerbe des équivalences \deleted{\unsure{en}} de la gerbe \deleted{$\mathscr G\unsure{(\mathbb Z_p)}$} \inserted{de gpe $\mu_2$}
définie sur $\mathbb Z[\frac1p]$ par $(C,T)$ \deleted{\unsure{avec la}} avec la gerbe
$g'_{A''}$ déjà envisagée plus haut (sur $A'$ c'est \unsure{celle|celles}
des \unsure{$\sqrt{\ }$} de $A'(1)$, sur $A_\infty$ c'est $g'_\infty$ déjà envisagée
plus haut) et qui sur $A_p$ est \uuline{la} \uline{gerbe} \deleted{des}\inserted{\uline{triviale}}.

% B5 | text | small middle-lower writing, then a boxed deletion
Cette gerbe est \deleted{\unsure{la}} triviale, et l'ens des \deleted{classes}\inserted{sections} à isom
\deleted{\unsure{équiv. des}} près est \deleted{le} torseur $\mathcal T$ sous $H^1(\mathbb A,\mu_2)$ déjà
\inserted{envisagé plus haut, qui est \unsure{ess.} indép. de $\xi$.}
\deleted{définie par les équiv \unsure{$(-)$} dit, \unsure{on aurait} une
\uline{section canonique} \deleted{\unsure{$\mathscr G$}} \unsure{$s(\xi)$} de $\mathscr G(\xi)$}

% B6 | text | far left opposite large cancellation
\margin{\deleted{\unsure{D'autre part la}}}

% B7 | text | bottom paragraph
Tout ceci posé, la théorie \inserted{de Deuring} des courbes elliptiques
d'inv. de Hasse nul nous fournit

% B8 | text | bottom item a
a) un élément $\xi_p\in\mathcal T_p$ (l'\uline{élément de Deuring})

% B9 | text | bottom item b
b) une section \deleted{correspondante} $\unsure{\bar\xi|\xi}$ de $\mathscr G(\xi_p)$, définie
à isom près, unique \uline{au signe près à l'infini},
dont la valeur en $p$ \deleted{\unsure{a comme}} est la section unité
\deleted{\unsure{can.} \unsure{unique}} de $\mathscr G(\xi_p)_{\mathbb Z_p}$, et la valeur en $\infty$ est la
polarisation donnée par $T$\ldots

% B10 | text | foot, concluding lines
On trouve donc une \deleted{section}\inserted{\unsure{élémt|élément}} $\{\xi\}$ de $\mathcal T$, dont la
comp. suivant $p$ est $\xi_p$.
% archive sheet number: 49
